\documentclass[12pt]{article}
\usepackage[utf8]{inputenc}
\usepackage[T1]{fontenc}
\usepackage{hyperref}
\title{Tai Hongjin Language Resources for K-12 Curriculum Development}
\author{Compiled by Claude Sonnet 5 (claude-sonnet-5) for Andr\'as Kornai}
\date{}
\begin{document}
\maketitle
\section{Introduction}
Tai Hongjin (Tai Ya, ISO 639-3 tiz) is a Southwestern Tai language of south-central Yunnan, with a small community in northern Thailand. Speaker estimates range from about 50,000 to 150,000. It is distinct from Tai Lue. The searches found only general descriptions.

\section{Primary Educational Resources}
None found. The language has no traditional script.

\section{Linguistic and Reference Materials}
\begin{itemize}
\item Wikipedia on Tai Ya, including its five dialects: \url{https://en.wikipedia.org/wiki/Tai_Ya_language}
\item Omniglot on Tai Hongjin, with a Latin romanization: \url{https://www.omniglot.com/writing/taihongjin.htm}
\end{itemize}

\section{Religious and Community Resources}
\begin{itemize}
\item Joshua Project profile of the Hongjin Tai in China: \url{https://joshuaproject.net/people_groups/18500/CH}
\item Asia Harvest people-group profile: \url{https://people-groups.asiaharvest.org/people-group-profiles/yunnan/hongjin-tai.pdf}
\end{itemize}
No scripture or community-run language organisation was confirmed.

\section{Digital and Audio Resources}
None found.

\section{Curriculum Development Context}
Little exists for this language. A standard orthography, a lexicon and teaching materials would have to be developed from scratch, and the dialect differences among the five varieties would need to be dealt with.
\end{document}